\documentclass{beamer}

% Theme selection
\usetheme{Madrid}
\usecolortheme{default}

% Packages
\usepackage[utf8]{inputenc}
\usepackage[T1]{fontenc}
\usepackage{amsmath}
\usepackage{amssymb}
\usepackage{graphicx}

% Title information
\title[Strategi Investasi Kompetitif]{Strategi Investasi Kompetitif: Pendekatan Game Theory dan Portofolio Log-Optimal}
\subtitle{Analisis Matematis Berdasarkan Permainan $\phi$ (\textit{The $\phi$-Game})}
\date{\today}

\begin{document}

% Slide 1: Title
\begin{frame}
    \titlepage
\end{frame}

% Slide 2: Daftar Isi
\begin{frame}{Daftar Isi}
    \tableofcontents
\end{frame}

\section{Pendahuluan}

% Slide 2: Pendahuluan - Model Pasar Saham
\begin{frame}{Pendahuluan: Model Pasar Saham}
    \begin{itemize}
        \item \textbf{Vektor Pasar ($X$):} 
        $X = (X_1, \dots, X_m)$ di mana $X_i \ge 0$ adalah \textit{return} (nilai kembalian) per unit modal.
        \item \textbf{Vektor Portofolio ($b$):} 
        $b = (b_1, \dots, b_m)$ dengan kendala:
        \begin{itemize}
            \item $b_i \ge 0$ (tanpa \textit{short selling}).
            \item $\sum_{i=1}^m b_i = 1$ (modal penuh).
        \end{itemize}
        \item \textbf{Modal Akhir ($S$):} 
        Variabel acak $S = b^t X = \sum b_i X_i$.
        \item \textbf{Konteks Stokastik:} 
        $X$ ditarik dari fungsi distribusi bersama $F(x)$ yang diasumsikan diketahui.
    \end{itemize}
\end{frame}

% Slide 3: Kerangka Game Theory (Zero-Sum)
\begin{frame}{Kerangka Game Theory (\textit{Zero-Sum})}
    \begin{block}{Definisi Permainan}
        \begin{itemize}
            \item \textbf{Pemain:} Dua investor (Pemain 1 vs Pemain 2).
            \item \textbf{Tujuan:} Memaksimalkan ekspektasi fungsi imbal hasil:
            $$E\phi(S_1/S_2)$$
            \item \textbf{Kriteria:} 
            \begin{itemize}
                \item Kemenangan ditentukan oleh rasio kekayaan $S_1/S_2$.
                \item $\phi(t)$ adalah fungsi monoton naik (\textit{non-decreasing}).
            \end{itemize}
        \end{itemize}
    \end{block}
    \begin{alertblock}{Hipotesis Utama}
        Keberadaan strategi universal \textbf{Portofolio Log-Optimal ($b^*$)} yang tangguh (\textit{robust}) terhadap berbagai fungsi $\phi$.
    \end{alertblock}
\end{frame}

\section{Kerangka Teoretis}

% Slide 4: Formalisasi Strategi Permainan
\begin{frame}{Formalisasi Strategi Permainan}
    \textbf{Ruang Strategi Dua Tahap:}
    \begin{enumerate}
        \item \textbf{Randomisasi Awal ($W$):} Memilih variabel acak adil (\textit{Fair Random Variable}) dengan $W \ge 0$ dan $E[W] \le 1$.
        \item \textbf{Alokasi Aset ($b$):} Memilih portofolio saham.
    \end{enumerate}
    
    \vspace{0.5cm}
    \textbf{Persamaan Payoff (1.1):}
    $$E \phi\left(\frac{W_{1}b_{1}^{\prime}X}{W_{2}b_{2}^{\prime}X}\right) = \iiint \phi\left(\frac{w_{1}b_{1}^{\prime}x}{w_{2}b_{2}^{\prime}x}\right) dG_{1} dG_{2} dF$$
    
    \vspace{0.5cm}
    \textbf{Masalah Minimax:}
    Mencari titik pelana ($v$) untuk keseimbangan strategi antara kedua pemain.
\end{frame}

% Slide 5: Permainan Primitif (phi-Game)
\begin{frame}{Permainan Primitif ($\phi$-Game)}
    \begin{itemize}
        \item \textbf{Definisi:} Permainan murni pemilihan distribusi probabilitas tanpa melibatkan pasar saham.
        \item \textbf{Masalah:} Memilih $W_1, W_2$ untuk memaksimumkan $E\phi(W_1/W_2)$.
        \item \textbf{Pertanyaan Kunci:}
        \begin{center}
            \textit{Kapan strategi optimal adalah strategi murni ($W=1$, tanpa judi)?}
        \end{center}
    \end{itemize}
\end{frame}

\section{Teorema Fundamental}

% Slide 6: Teorema 1 - Syarat Strategi Murni
\begin{frame}{Teorema 1: Syarat Strategi Murni}
    \begin{theorem}
        Permainan memiliki solusi optimal murni $W_{1}^{*} = W_{2}^{*} = 1$ jika dan hanya jika $\phi'(1) \ge 0$ dan untuk semua $t > 0$:
        $$\left(\frac{t-1}{t}\right)\phi'(1) \le \phi(t)-\phi(1) \le (t-1)\phi'(1)$$
    \end{theorem}
    
    \begin{itemize}
        \item \textbf{Interpretasi Geometris:} Fungsi $\phi$ harus bersifat \textbf{konkaf (cekung)} di sekitar titik 1.
        \item \textbf{Implikasi:} Untuk fungsi utilitas yang menghindari risiko (seperti $\ln t$ atau $t^\alpha$ dengan $0 \le \alpha \le 1$), strategi terbaik adalah "diam" (tidak berjudi).
    \end{itemize}
\end{frame}

% Slide 7: Keluarga Konveks (Convex Families)
\begin{frame}{Keluarga Konveks (\textit{Convex Families})}
    \begin{block}{Definisi}
        Himpunan variabel acak $\mathcal{S}$ disebut konveks jika untuk setiap $S_1, S_2 \in \mathcal{S}$ dan $0 \le \lambda \le 1$, berlaku:
        $$\lambda S_1 + (1-\lambda)S_2 \in \mathcal{S}$$
    \end{block}
    
    \textbf{Contoh Relevan:}
    \begin{itemize}
        \item Portofolio standar.
        \item Portofolio dengan kendala linear (misal: tanpa utang).
        \item Strategi sekuensial (multi-tahap).
    \end{itemize}
    
    \textit{Relevansi: Sifat ini memungkinkan penggunaan teorema minimax.}
\end{frame}

% Slide 8: Teorema 2 - Ekuivalensi Fundamental
\begin{frame}{Teorema 2: Ekuivalensi Logaritmik dan Linear}
    Untuk keluarga konveks $\mathcal{S}$, pernyataan berikut adalah ekuivalen untuk $S^* \in \mathcal{S}$:
    
    \begin{enumerate}
        \item \textbf{Sifat Log-Optimal:}
        $$E \ln(S/S^*) \le 0 \quad \text{untuk semua } S \in \mathcal{S}$$
        \item \textbf{Sifat Dominasi Linear (Supermartingale):}
        $$E(S/S^*) \le 1 \quad \text{untuk semua } S \in \mathcal{S}$$
    \end{enumerate}
    
    \vspace{0.5cm}
    \textbf{Sketsa Bukti:}
    \begin{itemize}
        \item Linear $\to$ Log: Menggunakan Ketidaksamaan Jensen.
        \item Log $\to$ Linear: Menggunakan Metode Variasi (perturbasi $\lambda \to 0$) dan kontradiksi.
    \end{itemize}
\end{frame}

% Slide 9: Teorema 3 - Pemisahan Variabel
\begin{frame}{Teorema 3: Pemisahan Variabel (\textit{Separation Theorem})}
    \begin{theorem}
        Strategi optimal dalam permainan pasar saham terfaktor menjadi dua komponen independen:
    \end{theorem}
    
    \begin{enumerate}
        \item \textbf{Komponen Judi ($W^*$):} 
        Solusi dari permainan primitif (bergantung pada preferensi risiko $\phi$).
        \item \textbf{Komponen Investasi ($b^*$):} 
        Selalu \textbf{Portofolio Log-Optimal}. Strategi ini bersifat universal dan tidak bergantung pada $\phi$.
    \end{enumerate}
    
    \vspace{0.3cm}
    \textbf{Mekanisme Bukti:}
    Menggunakan substitusi variabel $Z = \frac{W_2 S_2}{S^*}$ yang mereduksi masalah kompleks pasar saham kembali menjadi masalah permainan primitif sederhana.
\end{frame}

\section{Analisis Multi-tahap dan Stabilitas}

% Slide 10: Permainan Multi-tahap (Multistage)
\begin{frame}{Permainan Multi-tahap (\textit{Multistage})}
    \textbf{Model:}
    Investasi berulang selama $n$ periode dengan \textit{compounding} (bunga-berbunga) dan \textit{rebalancing}.
    
    \begin{block}{Teorema 4 (Strategi Myopic)}
        Strategi terbaik adalah memaksimalkan $E[\ln S]$ pada tiap langkah secara independen (\textit{greedy}).
    \end{block}
    
    \begin{itemize}
        \item \textbf{Dasar Matematis:} Sifat aditif fungsi logaritma:
        $$\ln\left(\prod_{k=1}^n b_k' X_k\right) = \sum_{k=1}^n \ln(b_k' X_k)$$
        \item \textbf{Implikasi:} Tidak perlu mengorbankan keuntungan hari ini demi strategi masa depan yang kompleks.
    \end{itemize}
\end{frame}

% Slide 11: Randomisasi Posterior (Observasi Modal)
\begin{frame}{Randomisasi Posterior: Mengintip Modal Lawan}
    \textbf{Skenario Baru:}
    Pemain dapat mengamati rasio modal $S_1/S_2$ di tengah permainan.
    
    \vspace{0.5cm}
    \textbf{Pertanyaan Psikologis:}
    \textit{Jika saya tertinggal jauh dari lawan, apakah saya harus mengubah strategi menjadi lebih agresif (Jockeying for position)?}
    
    \vspace{0.5cm}
    \textbf{Intuisi dari Judi Murni:}
    Dalam judi tanpa saham, pemain yang kalah memang seringkali harus mengambil risiko ekstrem ("All-in") untuk mengejar ketertinggalan.
\end{frame}

% Slide 12: Teorema 5 - Stabilitas Strategi
\begin{frame}{Teorema 5: Kemenangan Strategi "Buta"}
    \begin{alertblock}{Hasil Kontraintuitif}
        Di pasar saham, strategi optimal tetaplah \textbf{Log-Optimal ($b^*$)} dengan randomisasi awal standar ($W^*$).
    \end{alertblock}
    
    \vspace{0.3cm}
    \textbf{Mengapa demikian?}
    \begin{itemize}
        \item Portofolio $S^*$ memiliki sifat dominasi linear:
        $$E[S_{lawan}/S^*] \le 1$$
        \item Secara statistik, segala penyimpangan lawan dari $S^*$ (misal: spekulasi berlebih) justru akan merugikan modal relatif mereka sendiri.
        \item Oleh karena itu, kita tidak perlu bermanuver taktis (mengubah strategi). Cukup disiplin pada rencana awal.
    \end{itemize}
\end{frame}

\section{Kesimpulan}

% Slide 13: Kesimpulan
\begin{frame}{Kesimpulan}
    \begin{enumerate}
        \item \textbf{Penyatuan Teori:} 
        Portofolio Log-Optimal adalah solusi universal untuk berbagai tujuan investasi ($\phi$) dalam kerangka kompetitif.
        
        \item \textbf{Konsep "Fair Reach":} 
        Hukum $E(S/S^*) \le 1$ menjamin bahwa tidak ada strategi lain yang bisa secara konsisten mengalahkan log-optimal dalam jangka panjang.
        
        \item \textbf{Implikasi Praktis:} 
        Investor dapat fokus memaksimalkan pertumbuhan logaritmik portofolio sendiri (\textit{self-growth}) tanpa perlu bereaksi berlebihan terhadap pergerakan modal pesaing.
    \end{enumerate}
\end{frame}

\end{document}
